% Generated draft: requires mathematical and source review.
\documentclass[11pt]{article}
\usepackage[margin=25mm]{geometry}
\usepackage{fontspec}
\setmainfont{Latin Modern Roman}
\usepackage{amsmath,amssymb,mathtools,unicode-math}
\setmathfont{Latin Modern Math}
\usepackage{enumitem,needspace,xurl,hyperref}
\setlength{\parindent}{0pt}
\setlength{\parskip}{5pt}
\setlength{\emergencystretch}{3em}
\newcommand{\N}{\mathbb{N}}
\newcommand{\Z}{\mathbb{Z}}
\newcommand{\Q}{\mathbb{Q}}
\newcommand{\R}{\mathbb{R}}
\newcommand{\C}{\mathbb{C}}
\newcommand{\F}{\mathbb{F}}
\newcommand{\A}{\mathbb{A}}
\newcommand{\PP}{\mathbb{P}}
\newcommand{\E}{\mathbb{E}}
\newcommand{\eps}{\varepsilon}
\newcommand{\dd}{\,\mathrm{d}}
\begin{document}
\section*{UnsolvedMath ID: 30005598}
\textbf{Source problem code:} catalogue entry
\subsection{Definitions and mathematical statement}
Let \(\Omega\subset\R^2\) be a nonempty bounded domain with smooth boundary. To specify the usual uniform-field realization, for \(\beta>0\) set
\[
A_\beta(x_1,x_2)=\frac{\beta}{2}(-x_2,x_1),\qquad \partial_1(A_\beta)_2-\partial_2(A_\beta)_1=\beta.
\]
Define the lowest magnetic Neumann eigenvalue by
\[
\lambda_1^N(\Omega,\beta)
=\inf_{0\ne u\in H^1(\Omega;\C)}
\frac{\displaystyle\int_\Omega|(-i\nabla-A_\beta)u|^2\,\dd x}
{\displaystyle\int_\Omega|u|^2\,\dd x}.
\]
This is the lowest eigenvalue of \((-i\nabla-A_\beta)^2\) with the natural magnetic Neumann boundary condition
\[
\nu\cdot(-i\nabla-A_\beta)u=0\quad\text{on }\partial\Omega,
\]
where \(\nu\) is the outward unit normal.

The de Gennes constant is
\[
\Theta_0=\inf_{\xi\in\R}\ \inf_{0\ne v\in\mathcal V}
\frac{\displaystyle\int_0^\infty\bigl(|v'(t)|^2+(t-\xi)^2|v(t)|^2\bigr)\,\dd t}
{\displaystyle\int_0^\infty|v(t)|^2\,\dd t},
\qquad
\mathcal V=\{v\in H^1((0,\infty);\C):tv\in L^2((0,\infty);\C)\}.
\]
The associated half-line operators have Neumann condition \(v'(0)=0\). The supplied excerpt gives \(\Theta_0\approx0.590106\).

\textbf{Question under the positive-field, full-plane-potential interpretation.} Is it true that
\[
\forall\,\Omega\subset\R^2\text{ nonempty, bounded, smooth domain},\quad
\forall\,\beta>0,\qquad
\lambda_1^N(\Omega,\beta)\leq\Theta_0\beta?
\]
The supplied record says that even the disk case is open; this is a supplied status assessment, not an independently verified assertion here.

\textbf{Unresolved formulation qualifications.} The source statement does not specify the sign of \(\beta\). Its displayed right-hand side requires interpreting \(\beta\) as a nonnegative field strength; for a negative signed field the right-hand side is negative whereas the Rayleigh quotient is nonnegative. The formulation above uses \(\beta>0\), consistently with the excerpt's strict comparison with \(\beta\), but the intended treatment of zero and signed fields requires confirmation. The potential above is an explicit standard interpretation, not a recovered source hypothesis. On a domain with holes, prescribing the curl alone does not specify circulations around the holes. Thus the record must be checked to determine whether the intended operator uses the restriction of a full-plane uniform-field potential or allows additional circulation data.

\textbf{Other supplied context.} The excerpt attributes \(\lambda_1^N(\Omega,\beta)<\beta\) for bounded smooth planar domains to ``Theorem 4. [1]''; the identity of reference [1] is absent. It also records a lower bound \(\lambda_1^N(\Omega,\beta)\geq C\beta\) when \(\beta\delta^2\geq1\), with \(C\) universal, but does not define \(\delta\) or supply the complete hypotheses. Its damaged asymptotic sentence appears to describe quadratic small-field and linear large-field behavior; the exact notation and qualifications cannot be recovered reliably from the supplied extraction. These contextual statements are retained as source qualifications and are not additional precisely formulated claims in this draft.
\subsection{Short English statement}
The proposed bound says that the lowest magnetic Neumann eigenvalue of every bounded smooth planar domain is at most the de Gennes constant times the positive magnetic field strength. The statement above specifies the standard uniform-field operator explicitly. The field-sign convention, possible circulation data on domains with holes, and damaged contextual claims require source review.
\subsection{Sources}
\begin{itemize}
\item Local UnsolvedMath registry, \texttt{data/unsolvedmath/problems.json}, permanent ID \(30005598\), exact source problem code \texttt{catalogue entry}, title ``de Gennes Bound for Magnetic Neumann Eigenvalues.'' The supplied record belongs to the Oberwolfach Reports Open Problems collection. No individual proposer is supplied; the proposed year is recorded as 2023.
\item Supplied primary citation: \emph{Geometric Spectral Theory (2024)}, Oberwolfach Reports, pp.\ 2017--2106, DOI 10.4171/owr/2023/36. Source URL and literature-source link, labeled ``Oberwolfach Report 36'': \url{https://doi.org/10.4171/owr/2023/36}. The supplied title year is retained despite the 2023 DOI component and proposed-year metadata; their relationship has not been checked.
\item Supplied literature source: ``Magnetic Laplacian on the disc,'' \url{https://doi.org/10.1016/j.jmaa.2025.129261}. No authors or further bibliographic details are supplied.
\item Literature assessment dated exactly \texttt{2026-08-22}: supplied dated curation, not independently verified. It labels the proposed bound open, says that the primary report describes even the disk case as open, and reports that no proof or counterexample resolving that exact claim was found in the supplied literature triage. The record itself requests independent verification of time-sensitive status and cited claims.
\item The supplied statement status is \texttt{corrected\_verified}, with the assertion that the statement was reconstructed and checked against the official Oberwolfach report. This is registry provenance, not independent verification performed for this draft. The excerpt's citation ``[1]'' supporting its Theorem 4 cannot be identified from the supplied record.
\end{itemize}
\end{document}
